% Section 3 — pure body LaTeX. One \section; macros from main.tex only.
\section{Supersymmetric gauge theory and BV--BRST quantization}\label{sec:gaugebv}

This section collects the classical and quantum input for the one-loop computation of Sections~\ref{sec:superspace}--\ref{sec:component}: the general $\mathcal N=1$ gauge--chiral superfield action in the chiral representation (Section~\ref{sec:gaugebv:action}), its vector-frame reformulation (Section~\ref{sec:gaugebv:frames}) and component form (Section~\ref{sec:gaugebv:components}), the $\mathcal N=1,2,4$ super-Yang--Mills specializations (Section~\ref{sec:gaugebv:extended}), and the Batalin--Vilkovisky (BV) quantization~\cite{BV} with the Fermi--Feynman supergraph rules and the dimensional-reduction regulator (Sections~\ref{sec:gaugebv:bv}--\ref{sec:gaugebv:dred}). Two global conventions govern everything below. First, the gauge coupling is absorbed into the connection, $A_M=A_M^AT_A$, $\mathcal D_M=\partial_M-iA_M$; the canonical normalization is restored by $A_M=gA_M^{\rm can}$ (similarly for $\lambda,\phi,\psi,\mathscr D,F$), so every non-topological density carries an overall factor $g^{-2}$: vertices are $g$-independent up to that common factor, while propagators carry $g^{2}$. The compact gauge algebra is $[T_A,T_B]=if_{AB}{}^CT_C$ with orthonormal Killing metric $\kappa_{AB}=\delta_{AB}$, and $f^{ABC}:=\delta^{CD}f_{AB}{}^D$ is totally antisymmetric. Second, in Euclidean signature every tilded field is independent of its untilded partner before a middle-dimensional integration contour is chosen; on the Lorentzian contour the tilde becomes the formal conjugate. The computation of this paper is Euclidean; Lorentzian ancestors are quoted where instructive.

\subsection{Gauge and chiral superfields in the chiral representation}\label{sec:gaugebv:action}

Chiral and antichiral matter superfields $\Phi^I$, $\widetilde\Phi_I$, in a representation $\mathcal R_{\rm mat}$ and its dual, obey $\bcD_{\dot\alpha}\Phi^I=0$, $\cD_\alpha\widetilde\Phi_I=0$. Gauge transformations are generated by a chiral parameter $\Lambda=\Lambda^AT_A$ and an independent antichiral $\widetilde\Lambda$ (on the Lorentzian contour $\widetilde\Lambda=\bar\Lambda$):
\begin{equation}
h:=e^{i\Lambda},\qquad \widetilde h:=e^{i\widetilde\Lambda},\qquad \Phi'=h\Phi,\qquad \widetilde\Phi'=\widetilde\Phi\,\widetilde h^{-1}.
\end{equation}
The real prepotential $\cV=\cV^AT_A$ enters only through the bridge $e^{\cV}$, which dresses the flat spinor derivatives to gauge-covariant ones:
\begin{equation}\label{eq:gaugebv:bridge}
(e^{\cV})'=\widetilde h\,e^{\cV}h^{-1},\qquad
\Gamma_\alpha:=e^{-\cV}\cD_\alpha e^{\cV},\qquad
\widetilde\Gamma_{\dot\alpha}:=e^{\cV}\bcD_{\dot\alpha}e^{-\cV},\qquad
\nabla_\alpha:=\cD_\alpha+\Gamma_\alpha,\qquad
\nabla'_\alpha=h\nabla_\alpha h^{-1}.
\end{equation}
The gaugino field-strength superfields are
\begin{equation}\label{eq:gaugebv:W}
\cW_\alpha:=-\frac18\,\bcD^2\bigl(e^{-\cV}\cD_\alpha e^{\cV}\bigr),\qquad
\bcW_{\dot\alpha}:=+\frac18\,\cD^2\bigl(e^{\cV}\bcD_{\dot\alpha}e^{-\cV}\bigr),
\end{equation}
chiral and antichiral respectively and covariant: $\bcD_{\dot\alpha}\cW_\beta=0$, $\cD_\alpha\bcW_{\dot\beta}=0$, $\cW'_\alpha=h\cW_\alpha h^{-1}$, $\bcW'_{\dot\alpha}=\widetilde h\bcW_{\dot\alpha}\widetilde h^{-1}$. Engineering dimensions are $[\Phi]=1$, $[\cW_\alpha]=3/2$, $[\cV]=0$.

With the ordered Berezin projectors of Section~\ref{sec:notation}, $[X]_F:=-\tfrac14\cD^2X|$, $[\widetilde X]_{\widetilde F}:=-\tfrac14\bcD^2\widetilde X|$, $[Y]_D:=\tfrac1{16}\cD^2\bcD^2Y|$ (bar derivatives acting first), the general local two-derivative action is determined by four independent data~\cite{WeinbergIII,GGRS}: a real gauged K\"ahler density $\mathscr K_g(\widetilde\Phi,\Phi,e^\cV)$ of dimension $2$, a chiral superpotential $W(\Phi)$ of dimension $3$, a symmetric chiral kinetic matrix $f_{AB}(\Phi)$ of dimension $0$, and a Fayet--Iliopoulos parameter $\xi_A$. The exact Euclidean action used throughout this paper is
\begin{equation}\label{eq:gaugebv:SE}
\begin{aligned}
S_E=-\int d^4x\,\Big\{&
[\mathscr K_g(\widetilde\Phi,\Phi,e^\cV)]_D
+[W(\Phi)]_F+[\widetilde W(\widetilde\Phi)]_{\widetilde F}\\
&+\frac14[f_{AB}(\Phi)\,\cW^{A\alpha}\cW^B_\alpha]_F
+\frac14[\widetilde f_{AB}(\widetilde\Phi)\,\bcW^A_{\dot\alpha}\bcW^{B\dot\alpha}]_{\widetilde F}
+\xi_A[\cV^A]_D\Big\},
\end{aligned}
\end{equation}
the Wick image of its Lorentzian ancestor $S_L$, with $\mathcal L_E=-\mathcal L_L|_{\rm Wick}$. Exact gauge invariance of~\eqref{eq:gaugebv:SE} is equivalent to
\begin{equation}\label{eq:gaugebv:invariance}
X_A^I\,\partial_IW=0,\qquad
X_C^I\partial_If_{AB}-f_{CA}{}^Df_{DB}-f_{CB}{}^Df_{AD}=0,\qquad
\xi_Af_{BC}{}^A=0,
\end{equation}
with the holomorphic Killing vectors $X_A^I:=i(T_A)^I{}_J\Phi^J$. The canonical slice of~\eqref{eq:gaugebv:SE} is
\begin{equation}\label{eq:gaugebv:canonical}
\mathscr K_g=\widetilde\Phi_I(e^\cV)^I{}_J\Phi^J,\qquad
f_{AB}=\Bigl(\frac1{g^2}-\frac{i\theta_{\rm YM}}{8\pi^2}\Bigr)\delta_{AB},\qquad
\widetilde f_{AB}=\Bigl(\frac1{g^2}+\frac{i\theta_{\rm YM}}{8\pi^2}\Bigr)\delta_{AB}.
\end{equation}

\subsection{Chiral versus vector representation}\label{sec:gaugebv:frames}

The chiral representation is one of three gauge frames related by bridge factorization~\cite{GGRS}. Two complexified bridges $\mathcal B,\widetilde{\mathcal B}$ with
\begin{equation}\label{eq:gaugebv:frames}
e^\cV=\widetilde{\mathcal B}\,\mathcal B,\qquad
\mathcal B'=k\mathcal Bh^{-1},\qquad
\widetilde{\mathcal B}'=\widetilde h\,\widetilde{\mathcal B}k^{-1}
\end{equation}
interpolate between the chiral frame, in which matter transforms by the complexified parameters $(h,\widetilde h)$, and the vector (real) frame, in which matter $\Phi^{\mathsf V}:=\mathcal B\Phi$ transforms by the real group element $k$ only, $\Phi^{\mathsf V}{}'=k\Phi^{\mathsf V}$, and chirality is covariant, $\bar\nabla^{\mathsf V}_{\dot\alpha}\Phi^{\mathsf V}=0$. In the vector frame all superspace derivatives carry connections; imposing the conventional constraints
\begin{equation}\label{eq:gaugebv:constraints}
\mathcal F^{\mathsf V}_{\alpha\beta}=0,\qquad
\mathcal F^{\mathsf V}_{\dot\alpha\dot\beta}=0,\qquad
\mathcal F^{\mathsf V}_{\alpha\dot\beta}=0
\end{equation}
solves them in terms of the bridges, $\nabla^{\mathsf V}_\alpha=\widetilde{\mathcal B}^{-1}\circ\cD_\alpha\circ\widetilde{\mathcal B}$, $\bar\nabla^{\mathsf V}_{\dot\alpha}=\mathcal B\circ\bcD_{\dot\alpha}\circ\mathcal B^{-1}$, with the vector derivative defined by $\{\nabla^{\mathsf V}_\alpha,\bar\nabla^{\mathsf V}_{\dot\beta}\}=-2(\sigma_E^m)_{\alpha\dot\beta}\mathcal D^{\mathsf V}_m$ in Euclidean signature. The two representations agree identically at the level of the action: the density identity
\begin{equation}\label{eq:gaugebv:densityid}
\widetilde\Phi^{\mathsf V}\Phi^{\mathsf V}=\widetilde\Phi\,e^\cV\Phi
\end{equation}
and its field-strength analog imply $S^{\mathsf V}=S^{\mathsf C}$ term by term, the value of~\eqref{eq:gaugebv:SE}. The symmetric-bridge gauge $\mathcal B=\widetilde{\mathcal B}=e^{\cV/2}$ makes the identification manifest.

\begin{remark}\label{rem:gaugebv:chiralchoice}
The perturbative computation of this paper is formulated entirely in the chiral representation. In this frame the gauge orbit is parametrized by the chiral--antichiral pair $(\Lambda,\widetilde\Lambda)$, the Faddeev--Popov ghosts are chiral superfields (Section~\ref{sec:gaugebv:bv}), the gauge functions $\mathcal F_\pm$ are chiral and antichiral, and every interaction vertex is an ordered word in the single real prepotential $\cV$ (Section~\ref{sec:gaugebv:ff}); by~\eqref{eq:gaugebv:densityid} no physics is lost. The vector frame enters only through the background-covariant derivatives used to define $\mathcal F_\pm$ and the background-field split.
\end{remark}

\subsection{Component reconstruction}\label{sec:gaugebv:components}

The Wess--Zumino gauge is the slice in which $\cV$ is at most quadratic; the normalized Euclidean representative is
\begin{equation}\label{eq:gaugebv:WZ}
\cV^{\rm WZ}=-2i\,\theta\sigma_E^m\bar\theta\,A_m
+2i\theta^2\bar\theta_{\dot\alpha}\widetilde\lambda^{\dot\alpha}
-2i\bar\theta^2\theta^\alpha\lambda_\alpha
+\theta^2\bar\theta^2\mathscr D,
\end{equation}
with $\mathscr D$ the auxiliary field (not to be confused with the spinor derivatives $\cD_\alpha,\bcD_{\dot\alpha}$), $(\cV^{\rm WZ})^3=0$ so that $e^{\pm\cV}=1\pm\cV+\tfrac12\cV^2$, and ordinary gauge transformations as the residual group. Matter component fields are the ordered projections
\begin{equation}\label{eq:gaugebv:matterproj}
\phi^I:=\Phi^I|,\qquad
\psi_\alpha^I:=\frac1{\sqrt2}\cD_\alpha\Phi^I|,\qquad
F^I:=-\frac14\cD^2\Phi^I|,
\end{equation}
together with the tilded partners; the vector-multiplet components are
\begin{equation}\label{eq:gaugebv:vectorproj}
\lambda_\alpha:=i\cW_\alpha|,\qquad
\widetilde\lambda_{\dot\alpha}:=-i\bcW_{\dot\alpha}|,\qquad
\mathscr D:=-\frac12\cD^\alpha\cW_\alpha|=+\frac12\bcD_{\dot\alpha}\bcW^{\dot\alpha}|,\qquad
\cD_{(\alpha}\cW_{\beta)}|=-i(\sigma_E^{mn})_{\alpha\beta}F_{mn},
\end{equation}
the equality of the two $\mathscr D$ projections being the superspace Bianchi identity. In the canonical slice~\eqref{eq:gaugebv:canonical} the Euclidean component density $\mathcal L_E=\mathcal L_{E,\rm matter}+\mathcal L_{E,W}+\mathcal L_{E,\rm gauge}$ reads
\begin{equation}\label{eq:gaugebv:LEcan}
\begin{aligned}
\mathcal L_{E,\rm can}={}&
(\mathcal D_m\widetilde\phi)_I(\mathcal D_m\phi)^I
+\widetilde\psi_{\dot\alpha,I}(\bar\sigma_E^m)^{\dot\alpha\alpha}(\mathcal D_m\psi_\alpha)^I
-\widetilde F_IF^I-\mu_A^{(0)}\mathscr D^A\\
&-i\sqrt2\,\bigl[\widetilde\phi_I(T_A)^I{}_J\lambda^{A\alpha}\psi_\alpha^J
-\widetilde\psi_{\dot\alpha,I}(T_A)^I{}_J\widetilde\lambda^{A\dot\alpha}\phi^J\bigr]\\
&-W_I(\phi)F^I+\frac12W_{IJ}(\phi)\,\psi^{I\alpha}\psi_\alpha^J
-\widetilde W^{,I}(\widetilde\phi)\widetilde F_I
+\frac12\widetilde W^{,IJ}(\widetilde\phi)\,\widetilde\psi_{\dot\alpha,I}\widetilde\psi^{\dot\alpha}_J\\
&+\frac1{4g^2}F^A_{mn}F^A_{mn}
+\frac1{g^2}\widetilde\lambda^A_{\dot\alpha}(\bar\sigma_E^m)^{\dot\alpha\alpha}(\mathcal D_m\lambda_\alpha)^A
-\frac1{2g^2}\mathscr D^A\mathscr D^A
+\frac{i\theta_{\rm YM}}{64\pi^2}\epsilon_E^{mnrs}F^A_{mn}F^A_{rs},
\end{aligned}
\end{equation}
with $\mu_A^{(0)}:=\widetilde\phi_I(T_A)^I{}_J\phi^J+\xi_A$ and $\epsilon_E^{1234}=+1$. The auxiliary fields enter algebraically; their equations of motion
\begin{equation}\label{eq:gaugebv:auxeom}
F^I=-\widetilde W^{,I}(\widetilde\phi),\qquad
\widetilde F_I=-W_I(\phi),\qquad
\mathscr D^A=-g^2\delta^{AB}\mu_B^{(0)}
\end{equation}
complete the squares of the positive bosonic contour.

\begin{remark}\label{rem:gaugebv:componentmap}
The complete off-shell component map --- including the general K\"ahler geometry of the target, the moment map, and the shifted auxiliaries --- has been reconstructed in the source repository and machine-verified coefficient by coefficient ($336$ independent component coefficients, zero failures); only the canonical slice is displayed here. The component-field route to the anomaly (Section~\ref{sec:component}) uses this map and is partial in the sense stated there.
\end{remark}

\subsection{$\mathcal N=1$, $\mathcal N=2$ and $\mathcal N=4$ super-Yang--Mills}\label{sec:gaugebv:extended}

Pure $\mathcal N=1$ super-Yang--Mills is the vector multiplet $(A_m^A,\lambda_\alpha^A,\widetilde\lambda_{\dot\alpha}^A,\mathscr D^A)$ with action
\begin{equation}\label{eq:gaugebv:SEN1}
S_E^{(1)}=-\frac14\int d^4x\,\bigl(
[f_{AB}\cW^{A\alpha}\cW^B_\alpha]_F
+[\widetilde f_{AB}\bcW^A_{\dot\alpha}\bcW^{B\dot\alpha}]_{\widetilde F}\bigr).
\end{equation}
$\mathcal N=2$ super-Yang--Mills adds one adjoint chiral multiplet $\Phi=(\phi,\psi_\alpha,F)$; closure of the second supersymmetry at the linearized level, together with gauge invariance, fixes the relative coefficient of the matter kinetic term uniquely, giving
\begin{equation}\label{eq:gaugebv:SEN2}
S_E^{(2)}=-\int d^4x\,\Big\{
\frac1{g^2}\delta_{AB}\bigl[\widetilde\Phi^A(e^{\cV_{\rm ad}})^B{}_C\Phi^C\bigr]_D
+\frac14[f_{AB}\cW^{A\alpha}\cW^B_\alpha]_F
+\frac14[\widetilde f_{AB}\bcW^A_{\dot\alpha}\bcW^{B\dot\alpha}]_{\widetilde F}\Big\},
\end{equation}
with $e^{\cV_{\rm ad}}$ the bridge in the adjoint representation. $\mathcal N=4$ super-Yang--Mills, the theory studied in this paper, contains three adjoint chirals $\Phi_r$, $r=1,2,3$, an $SU(3)$ flavor triplet inside $SU(4)_R$, together with the $SU(3)$-invariant superpotential
\begin{equation}\label{eq:gaugebv:WN4}
W=-\frac{\sqrt2}{6g^2}\,\varepsilon_{rst}f^{ABC}\,\Phi_r^A\Phi_s^B\Phi_t^C,
\end{equation}
whose coefficient $-\sqrt2/6$ is fixed uniquely by matching the gauge Yukawa coupling. The action is
\begin{equation}\label{eq:gaugebv:SEN4}
\begin{aligned}
S_E^{(4)}=-\int d^4x\,\Big\{&
\frac1{g^2}\delta_{AB}\sum_{r=1}^3\bigl[\widetilde\Phi_r^A(e^{\cV_{\rm ad}})^B{}_C\Phi_r^C\bigr]_D
+\frac14[f_{AB}\cW^{A\alpha}\cW^B_\alpha]_F
+\frac14[\widetilde f_{AB}\bcW^A_{\dot\alpha}\bcW^{B\dot\alpha}]_{\widetilde F}
+[W]_F+[\widetilde W]_{\widetilde F}\Big\}.
\end{aligned}
\end{equation}
The unified $\mathcal N=1,2,4$ form, with multiplet counts $m_1=0$, $m_2=1$, $m_4=3$, is the starting point of the vertex inventory of Section~\ref{sec:gaugebv:ff}.

\subsection{BV--BRST quantization}\label{sec:gaugebv:bv}

Because the even noncompact gauge-orbit integral diverges and the odd one vanishes, gauge fixing is mandatory, and we use the BV formalism~\cite{BV}. The minimal sector consists of the physical fields and a chiral--antichiral pair of Faddeev--Popov ghost superfields $c,\widetilde c$; the non-minimal sector consists of an antighost--Nakanishi--Lautrup doublet $(\bar c_+,b_+)$ of chiral superfields and its antichiral partner $(\bar c_-,b_-)$, with $s\bar c_\pm=b_\pm$, $sb_\pm=0$. Every field $X$ carries a Grassmann parity $\epsilon_X$, a ghost number and an engineering dimension; its BV antifield $X^*$ (never complex conjugation) obeys
\begin{equation}\label{eq:gaugebv:antifield}
\epsilon_{X^*}=\epsilon_X+1,\qquad
\operatorname{gh}(X^*)=-1-\operatorname{gh}(X),\qquad
[X^*]=d_\Sigma-[X],
\end{equation}
with $d_\Sigma=2$ for the full superspace and $d_\Sigma=3$ for the chiral and antichiral subspaces. The complete spectrum is
\begin{equation}\label{eq:gaugebv:spectrum}
\begin{array}{c|c|c|c|c}
X&\text{domain}&\epsilon_X&\operatorname{gh}(X)&[X]\\ \hline
\cV&(x,\theta,\bar\theta)&0&0&0\\
\Phi^I&(y,\theta)&0&0&1\\
\widetilde\Phi_I&(\widetilde y,\bar\theta)&0&0&1\\
c&(y,\theta)&1&+1&0\\
\widetilde c&(\widetilde y,\bar\theta)&1&+1&0\\
\bar c_\pm&(y,\theta),\,(\widetilde y,\bar\theta)&1&-1&2\\
b_\pm&(y,\theta),\,(\widetilde y,\bar\theta)&0&0&2
\end{array}
\end{equation}
The classical BRST differential $s$ ($\epsilon=1$, $\operatorname{gh}=1$) acts as
\begin{equation}\label{eq:gaugebv:brst}
se^\cV=i\widetilde c\,e^\cV-ie^\cV c,\qquad
sc=ic^2,\qquad
s\widetilde c=i\widetilde c^{\,2},\qquad
s\Phi=ic\Phi,\qquad
s\widetilde\Phi=-i\widetilde\Phi\,\widetilde c,
\end{equation}
in components $sc^C=-\tfrac12 f_{AB}{}^Cc^Ac^B$, and is nilpotent on the untruncated local field algebra: $s^2=0$ and $sS_0=0$. The exact transformation of the prepotential is the Bernoulli series
\begin{equation}\label{eq:gaugebv:brstV}
s\cV=\frac{\operatorname{ad}_\cV}{1-e^{-\operatorname{ad}_\cV}}\bigl(ie^{-\operatorname{ad}_\cV}\widetilde c-ic\bigr)
=i\sum_{p,q\ge0}\frac{B_{p+q}}{p!\,q!}\bigl[(-1)^q\cV^p\widetilde c\,\cV^q-(-1)^p\cV^p c\,\cV^q\bigr],
\end{equation}
with $B_0=1$, $B_1=-\tfrac12$, $B_2=\tfrac16$, $B_3=0$.

With the ordered pairings $\langle X,Y\rangle_8:=\int d^4x\,[X_AY^A]_D$ and their chiral/antichiral analogs $\langle\cdot,\cdot\rangle_\pm$, the antibracket and the minimal master action are
\begin{equation}\label{eq:gaugebv:antibracket}
(F,G):=\sum_X\int_{\Sigma_X}\Bigl[
F\frac{\overleftarrow\delta}{\delta X}\frac{\vec\delta G}{\delta X^*}
-F\frac{\overleftarrow\delta}{\delta X^*}\frac{\vec\delta G}{\delta X}\Bigr],
\end{equation}
\begin{equation}\label{eq:gaugebv:Smin}
S_{\min}:=S_0
+\langle\cV^*,s\cV\rangle_8
+\langle\Phi^*,s\Phi\rangle_+
+\langle\widetilde\Phi^*,s\widetilde\Phi\rangle_-
-\langle c^*,sc\rangle_+
-\langle\widetilde c^{\,*},s\widetilde c\rangle_-,
\end{equation}
supplemented by the non-minimal stabilization $S_{\rm nm}=S_{\min}+\langle\bar c^{\,*},b\rangle$. They obey the classical master equation and generate $s$ as a Hamiltonian vector field:
\begin{equation}\label{eq:gaugebv:CME}
\frac12(S_{\min},S_{\min})=0,\qquad
sF=(S_{\min},F),\qquad
s^2F=\frac12\bigl((S_{\min},S_{\min}),F\bigr)=0.
\end{equation}

For the background split $e^{\cV_{\rm tot}}=\widetilde{\mathcal B}_{\rm B}e^{\cV_{\rm q}}\mathcal B_{\rm B}$ the gauge functions are the background-covariant chiral and antichiral projections of the quantum prepotential,
\begin{equation}\label{eq:gaugebv:Fpm}
\mathcal F_+:=-\frac14(\bar\nabla^{\mathsf V}_{\rm B})^2\cV_{\rm q},\qquad
\bar\nabla^{\mathsf V}_{{\rm B}\dot\alpha}\mathcal F_+=0,\qquad
\mathcal F_-:=-\frac14(\nabla^{\mathsf V}_{\rm B})^2\cV_{\rm q},\qquad
\nabla^{\mathsf V}_{{\rm B}\alpha}\mathcal F_-=0,\qquad
[\mathcal F_\pm]=1.
\end{equation}
With the collective pairing of the chiral/antichiral sectors, the standard gauge fermion and its BRST variation are
\begin{equation}\label{eq:gaugebv:sPsi}
\Psi:=\langle\bar c,\mathcal F\rangle-\frac12\langle\bar c,\mathcal Yb\rangle,\qquad
s\Psi=\langle b,\mathcal F\rangle
-\langle\bar c,\mathcal M^{\rm FP}c^{\rm pair}\rangle
-\frac12\langle b,\mathcal Yb\rangle
+\frac12\langle\bar c,(s\mathcal Y)b\rangle,
\end{equation}
where $\mathcal Y$ is an even, invertible, local, background-covariant kernel of dimension $-1$, $c^{\rm pair}:=(c,\widetilde c)$ dressed to the background, and the Faddeev--Popov operator $\mathcal M^{\rm FP}$ (even, ghost-number zero, dimension $1$) is the gauge tangent of the quantum field, $s\mathcal F=\mathcal M^{\rm FP}c^{\rm pair}$. Integrating out the non-minimal doublets produces the Faddeev--Popov Berezinian
\begin{equation}\label{eq:gaugebv:FPfactor}
\mathfrak F^{\rm FP}=\operatorname{Ber}\bigl(\mathcal M_0^{-1}\mathcal M^{\rm FP}\bigr),
\end{equation}
evaluated on the quotient by the residual stabilizer, and, whenever the kernel $\mathcal Y$ is itself gauge averaged, a Nielsen--Kallosh factor fixed by the requirement that the multiplier Gaussian times it reproduce the declared gauge-average normalization. At vanishing external antifield sources the gauge-fixed action is $S_\Psi=S_0+s\Psi$.

The quantum master action $W=S_{\min}+\sum_{n\ge1}\hbar^nM_n$ must satisfy the quantum master equation
\begin{equation}\label{eq:gaugebv:QME}
\frac12(W,W)-i\hbar\Delta W=0\ \ (\text{Lorentzian}),\qquad
\frac12(W,W)-\hbar\Delta W=0\ \ (\text{Euclidean}),
\end{equation}
with $\Delta$ the BV Laplacian of the measure; a nonzero left-hand side is the regulated BV anomaly functional. Expanded in $\hbar$ this is a recursion whose first (one-loop) rung in the Euclidean theory reads $sM_1=\Delta S_{\min}$.

All Ward identities descend from the single Schwinger--Dyson identity of the regulated measure: for every coefficient mode $q$ of the cylindrical expansion, with measure density $\varpi$ and total path-integral exponent $\mathcal E_{\rm tot}$,
\begin{equation}\label{eq:gaugebv:SD}
0=\int d\mu\;e^{\mathcal E_{\rm tot}}\,
\Bigl[\varpi^{-1}\frac{\vec\partial\varpi}{\partial q}+\frac{\vec\partial\,\mathcal E_{\rm tot}}{\partial q}\Bigr].
\end{equation}
BRST, background and supersymmetry Ward identities are specializations of~\eqref{eq:gaugebv:SD} to the corresponding vector fields. The anomalous supersymmetry Ward identity of this paper is the statement that the supersymmetry specialization of~\eqref{eq:gaugebv:SD} retains a nonzero defect at one loop; it is constructed in Section~\ref{sec:superspace}.

\begin{remark}\label{rem:gaugebv:cutoffdefect}
At a finite mode cutoff the classical master equation is not automatic: the truncated master action obeys $\tfrac12(S_{\min,\nu},S_{\min,\nu})_\nu=s_\nu S_{0,\nu}+\sum(-1)^{\epsilon}Q^*\,s_\nu^2Q$, whose right-hand side vanishes only after regulator closure, regulator invariance of the action, and nilpotency of the truncated differential are established. The identities above are therefore always understood with their regulated defect terms retained.
\end{remark}

\subsection{Gauge fixing and propagators}\label{sec:gaugebv:ff}

We now fix the perturbative scheme used in Sections~\ref{sec:superspace}--\ref{sec:heatkernel}. The Fermi--Feynman representative is the kernel choice
\begin{equation}\label{eq:gaugebv:YFF}
\mathcal Y_{\rm FF}^{-1}=\frac1{2g^2}\,\mathcal A,\qquad
\mathcal A:=\begin{pmatrix}0&-\bcD^2/4\\[2pt]-\cD^2/4&0\end{pmatrix},\qquad
\mathcal A^2=\BoxE\,\mathbf1,\qquad
\mathcal Y_{\rm FF}=2g^2\,\frac{\mathcal A}{\BoxE},
\end{equation}
acting on the gauge-function doublet $(\mathcal F_+,\mathcal F_-)$; it produces the gauge-fixing term $\tfrac12\langle\mathcal Y_{\rm FF}^{-1}\mathcal F,\mathcal F\rangle$.

\begin{remark}\label{rem:gaugebv:ffscheme}
Two scheme choices are declared here. First, $\mathcal Y_{\rm FF}$ contains $\BoxE^{-1}$ and is therefore nonlocal as a gauge-fermion kernel: perturbation theory is \emph{defined} by Fermi--Feynman Gaussian averaging rather than derived from a local admissible gauge fermion. (The source repository explicitly records this locality obstruction and does not assert a unique propagator set, a selected Nielsen--Kallosh branch, or momentum-space rules; the rules below are the standard Fermi--Feynman supergraph grammar~\cite{GGRS}, with coefficients that agree with the verified repo expressions.) Second, because $\mathcal Y_{\rm FF}$ is field independent, $s\mathcal Y_{\rm FF}=0$ and the associated Nielsen--Kallosh factor is a field-independent constant in this slice; it cancels in normalized correlation functions and is set to unity.
\end{remark}

With $\BoxE:=\delta^{mn}\partial_m\partial_n$ and $z=(x,\theta,\bar\theta)$, the chiral, antichiral and transverse projectors are
\begin{equation}\label{eq:gaugebv:projectors}
\mathcal P_+:=\frac{\bcD^2\cD^2}{16\BoxE},\qquad
\mathcal P_-:=\frac{\cD^2\bcD^2}{16\BoxE},\qquad
\mathcal P_T:=-\frac{\cD^\alpha\bcD^2\cD_\alpha}{8\BoxE},\qquad
\mathcal P_0:=\mathcal P_++\mathcal P_-,
\end{equation}
with the projector algebra following from the spinor-derivative identities:
\begin{equation}\label{eq:gaugebv:dalgebra}
\mathcal P_i\mathcal P_j=\delta_{ij}\mathcal P_i\ \ (i,j\in\{T,+,-\}),\qquad
\mathcal P_T+\mathcal P_++\mathcal P_-=1,
\end{equation}
\begin{equation}\label{eq:gaugebv:Didentities}
\bcD^2\cD^2\bcD^2=16\BoxE\,\bcD^2,\qquad
\cD^2\bcD^2\cD^2=16\BoxE\,\cD^2,\qquad
-\cD^\alpha\bcD^2\cD_\alpha=8\BoxE-\frac12\{\cD^2,\bcD^2\}.
\end{equation}
The Fermi--Feynman quadratic form $S^{(2)}_{E,V}=\tfrac1{4g^2}\int d^4x\,d^4\theta\;\cV^A\BoxE\cV^A$ and the constrained identity kernels $\mathbf1_\pm:=\mathcal P_\pm\delta^8(z-z')$ give the exact vector and matter propagators
\begin{equation}\label{eq:gaugebv:propagators}
G^{VV,AB}=2\h g^2\delta^{AB}\,\BoxE^{-1},\qquad
G^{\Phi\widetilde\Phi,AB}=-\h g^2\delta^{AB}\,\mathcal P_+\,\delta^8(z-z'),
\end{equation}
the reversed matter orientation using $\mathcal P_-$. In momentum space (phase $e^{ipx}$, all momenta incoming, $\BoxE\mapsto-p^2$; $\theta_{12}:=\theta_1-\theta_2$):
\begin{equation}\label{eq:gaugebv:momprop}
\begin{aligned}
\langle\cV^A(p,\theta_1)\cV^B(p',\theta_2)\rangle
&=-(2\pi)^4\delta^4(p+p')\,\frac{2\h g^2\delta^{AB}}{p^2}\,\delta^4(\theta_{12}),\\
\langle\Phi^A(p,\theta_1)\widetilde\Phi^B(p',\theta_2)\rangle
&=(2\pi)^4\delta^4(p+p')\,\frac{\h g^2\delta^{AB}}{16p^2}\,\bcD_1^2\cD_1^2\delta^4(\theta_{12}).
\end{aligned}
\end{equation}
The spinor-derivative algebra collapses products of matter propagators at coincident superspace points through the loop-saturation identity
\begin{equation}\label{eq:gaugebv:saturation}
\delta^4(\theta_{12})\,\cD^2\bcD^2\,\delta^4(\theta_{12})=16\,\delta^4(\theta_{12}).
\end{equation}

The interaction vertices are ordered superfield words; we abbreviate $\int_8:=\int d^4x\,d^4\theta$ and $\int_\pm$ for the chiral/antichiral integrals. Expanding~\eqref{eq:gaugebv:W} with the exact exponential derivatives
\begin{equation}\label{eq:gaugebv:expderivative}
e^{-\cV}\cD_\alpha e^{\cV}=\sum_{p,q\ge0}\frac{(-1)^p}{p!\,q!\,(p+q+1)}\,\cV^p(\cD_\alpha\cV)\cV^q,\qquad
e^{\cV}\bcD_{\dot\alpha}e^{-\cV}=\sum_{p,q\ge0}\frac{(-1)^{q+1}}{p!\,q!\,(p+q+1)}\,\cV^p(\bcD_{\dot\alpha}\cV)\cV^q,
\end{equation}
the pure-gauge words are
\begin{equation}\label{eq:gaugebv:gaugevertex}
S^{\rm g}_{E,+}=-\sum_{p,q,r,s\ge0}
\frac{(-1)^{p+r}f_{AB}}{256\,p!q!r!s!(p+q+1)(r+s+1)}
\int_+\bigl[\bcD^2(\cV^p\cD^\alpha\cV\,\cV^q)\bigr]^A
\bigl[\bcD^2(\cV^r\cD_\alpha\cV\,\cV^s)\bigr]^B,
\end{equation}
the antichiral words carrying $(-1)^{q+s}\widetilde f_{AB}$ with $\cD\leftrightarrow\bcD$; at cubic order this gives the four words with coefficients $\pm f_{AB}/512$. The matter--gauge chain is
\begin{equation}\label{eq:gaugebv:matchain}
S^{\rm mat}_E=-\frac1{g^2}\delta_{AB}\sum_{r=1}^{m_N}\sum_{n\ge0}\frac1{n!}
\int_8\widetilde\Phi_r^A\,\cV^{A_1}\cdots\cV^{A_n}(T_{A_1}\cdots T_{A_n})^B{}_C\,\Phi_r^C,
\end{equation}
with $(T_{A_1}\cdots T_{A_n})^B{}_C=i^nf_{A_1D_1}{}^Bf_{A_2D_2}{}^{D_1}\cdots f_{A_nC}{}^{D_{n-1}}$ in the adjoint representation; the labeled-leg vertex carries the color word symmetrized over its $n$ legs, and the cubic vertex is $-\tfrac{i}{g^2}f_{ABC}\int_8\widetilde\Phi^A\cV^B\Phi^C$. The $\mathcal N=4$ superpotential contributes the chiral and antichiral cubic primitives
\begin{equation}\label{eq:gaugebv:Wcubic}
\mathcal V_{\Phi\Phi\Phi}=\mathcal V_{\widetilde\Phi\widetilde\Phi\widetilde\Phi}=\frac{\sqrt2}{g^2}\,\varepsilon_{rst}f^{ABC}.
\end{equation}
The ghost sector is generated by inserting~\eqref{eq:gaugebv:brstV} into
\begin{equation}\label{eq:gaugebv:ghost}
S^{\rm FP}=\frac14\int_+\bar c_+\,\bcD^2(s\cV)+\frac14\int_-\bar c_-\,\cD^2(s\cV),\qquad
S^{(2)}_{\rm FP}=\frac i4\int_+\bar c_+\,\bcD^2\widetilde c-\frac i4\int_-\bar c_-\,\cD^2c,
\end{equation}
with ghost--ghost--vector words of coefficients $\tfrac i4\tfrac{B_{p+q}(-1)^q}{p!q!}$ (for $\widetilde c$) and $-\tfrac i4\tfrac{B_{p+q}(-1)^p}{p!q!}$ (for $c$), the lowest being the $B_1=-\tfrac12$ word $s\cV|_{(1)}=-\tfrac i2[\cV,c+\widetilde c]$. Supergraph weights follow the standard rule: a factor $\tau$ per vertex and $-\tau^{-1}$ per propagator, with $\tau=i/\hbar$ Lorentzian and $\tau=-1/\hbar$ Euclidean, graded permutation signs for labeled legs, and division by the automorphism group of the labeled graph.

\subsection{Regularization: dimensional reduction}\label{sec:gaugebv:dred}

Loop integrations are regularized by dimensional reduction~\cite{SiegelDRED}: loop momenta are $d$-dimensional, $d=4-2\epsilon$, while the spinor algebra, the superspace derivatives and all numerator spin words remain four-dimensional. Two metric representations coexist and must not be confused. The loop-momentum split is
\begin{equation}\label{eq:gaugebv:dredloop}
\delta_4^{mn}=\widehat\delta^{mn}+\breve\delta^{mn},\qquad
\widehat\delta^m{}_m=d,\qquad
\breve\delta^m{}_m=4-d=2\epsilon,\qquad
\widehat\delta^m{}_r\breve\delta^r{}_n=0,
\end{equation}
with purely hatted loop momenta, $\ell^m=\widehat\delta^m{}_n\ell^n$; the numerator spinor/vector split is
\begin{equation}\label{eq:gaugebv:dredspin}
\delta_d^{MN}=\bar\delta^{MN}+\widetilde\delta^{MN},\qquad
\operatorname{tr}\bar\delta=4,\qquad
\operatorname{tr}\widetilde\delta=d-4=-2\epsilon.
\end{equation}
The regulator enters the computation through a single axiom --- finite superspace spin words are contracted with the four-dimensional metric $\bar\delta$, while propagator inverses use the $d$-dimensional metric $\delta_d$ --- and through the evanescent scalar
\begin{equation}\label{eq:gaugebv:mul}
\mul^2:=\bar\ell^2-\ell_d^2=-\widetilde\ell^{\,2},
\end{equation}
the squared length of the loop momentum along the $-2\epsilon$ evanescent directions. Fourier and measure conventions are
\begin{equation}\label{eq:gaugebv:measure}
f(x)=\int\frac{d^4p}{(2\pi)^4}\,e^{ipx}f(p),\qquad
\BoxE\mapsto-p^2,\qquad
\int_\ell:=\mu^{2\epsilon}\int\frac{d^d\ell}{(2\pi)^d},
\end{equation}
with all external momenta taken incoming and the superspace delta functions $\delta^4(\theta_{12})$ of Section~\ref{sec:gaugebv:ff}.

\begin{remark}\label{rem:gaugebv:dredstatus}
Dimensional reduction is the only scheme-sensitive input of the one-loop computation: the single axiom above, together with the master integrals quoted in Section~\ref{sec:component}, is the complete regulator data used. Whether evanescent operators are required beyond this axiom is a question about higher-order structure and is deferred to the discussion of Section~\ref{sec:superspace}.
\end{remark}
